% header.tex —— research-to-book 导言区
% 由 template.tex 在 \begin{document} 之前 \input

% =========================================================
% 字体（XeLaTeX + ctexbook 已经处理 CJK 主字体，这里做精细分工）
% =========================================================
\usepackage{fontspec}
% 西文主字体：Times New Roman（macOS 自带）
\IfFontExistsTF{Times New Roman}{%
  \setmainfont{Times New Roman}%
}{%
  \setmainfont{Latin Modern Roman}%
}
\IfFontExistsTF{Helvetica Neue}{%
  \setsansfont{Helvetica Neue}%
}{%
  \setsansfont{Latin Modern Sans}%
}
\IfFontExistsTF{Menlo}{%
  \setmonofont{Menlo}[Scale=0.88]%
}{%
  \setmonofont{Latin Modern Mono}%
}

% CJK 字体分工
% 关键：macOS 的 STHeiti.ttc 在 XeTeX 下 glyph 映射会错位（已知 bug），
% 必须彻底替换掉 ctex 默认的 \heiti 绑定。
\usepackage{xeCJK}
% 主字体：宋体（正文）。不指定 BoldFont，让 xeCJK 自动用"加粗"（合成粗）
% —— 指定 BoldFont 找不到时会回退到 STHeiti，那是 bug 源头。
\setCJKmainfont{Songti SC}[AutoFakeBold=2.5]
\setCJKsansfont{PingFang SC}[AutoFakeBold=2.0]   % 无衬线
\setCJKmonofont{Songti SC}[AutoFakeBold=2.5]     % 等宽 CJK（代码块内中文）
% 关键：把 ctex 的 \heiti / \songti / \fangsong / \kaishu 全部重绑
\setCJKfamilyfont{heiti}{PingFang SC}[AutoFakeBold=2.0]   % 黑体 → PingFang（避开 STHeiti bug）
\setCJKfamilyfont{song}{Songti SC}[AutoFakeBold=2.5]      % 宋体
\setCJKfamilyfont{kaiti}{STKaiti}                         % 楷体（引文框）
\setCJKfamilyfont{fs}{STFangsong}                         % 仿宋
% 显式覆盖 ctex 的 \heiti 定义，确保万无一失
\renewcommand{\heiti}{\CJKfamily{heiti}}
\newcommand{\kai}{\CJKfamily{kaiti}}

% =========================================================
% 版式
% =========================================================
\usepackage[
  papersize={176mm,250mm},
  margin=22mm,
  top=25mm, bottom=25mm,
  headheight=14pt, headsep=8mm, footskip=10mm
]{geometry}

\usepackage{setspace}
\setstretch{1.45}

% =========================================================
% 颜色（朱砂 + 墨色 + 灰）
% =========================================================
\usepackage{xcolor}
\definecolor{zhusha}{HTML}{9B2D20}     % 朱砂红
\definecolor{mohei}{HTML}{1A1A1A}      % 墨黑
\definecolor{qianhui}{HTML}{F5F0EC}    % 浅米灰（引文框背景）
\definecolor{shenhui}{HTML}{555555}    % 深灰

% =========================================================
% 工具宏包
% =========================================================
\usepackage{booktabs}
\usepackage{longtable}
\usepackage{array}
\usepackage{tabularx}
% 数学公式（第 5、6 章公式密集；build.py 数学直通配合 Pandoc 输出 $...$/$$...$$）
\usepackage{amsmath}
\usepackage{amssymb}
% 图表标题：图/表名 + 按章编号 + 小号字
% 注意：不能用 caption 宏包——它与 pandoc longtable 输出的 \LTcaptype{none} 冲突
% （caption 包把 none 当计数器名，报 "No counter 'none' defined"）
\renewcommand{\figurename}{图}
\renewcommand{\tablename}{表}
\counterwithin{figure}{chapter}
\counterwithin{table}{chapter}
\makeatletter
\long\def\@makecaption#1#2{%
  {\small \heiti #1}\hspace{1em}{\small #2\par}}
\makeatother
\usepackage{calc}      % Pandoc 3.x 表格列宽计算需要
\usepackage{xfp}       % \real 命令
\usepackage[most]{tcolorbox}
% 注意：不加载 titlesec / titletoc，它们与 ctexset 冲突，会让目录折行。
% 全部章节/目录样式用 ctexset 统一控制。
\usepackage{fancyhdr}
\usepackage{enumitem}
\setlist{nosep, leftmargin=2em}
\providecommand{\tightlist}{\setlength{\itemsep}{0pt}\setlength{\parskip}{0pt}}
\usepackage{graphicx}
% pandoc 3.x 输出图片用的包装宏：图片宽于版心时等比缩到版心宽
\makeatletter
\newcommand{\pandocbounded}[1]{%
  \sbox\@tempboxa{#1}%
  \ifdim\wd\@tempboxa>\linewidth
    \resizebox{\linewidth}{!}{\usebox\@tempboxa}%
  \else
    \usebox\@tempboxa%
  \fi
}
\makeatother
\usepackage{microtype}
\usepackage[
  colorlinks=true,
  linkcolor=zhusha,
  urlcolor=zhusha,
  citecolor=zhusha,
  bookmarksnumbered=true,
  pdfstartview=Fit,
  pdftitle={{{BOOK_TITLE}}},
  pdfauthor={李佳伦}
]{hyperref}
\usepackage{bookmark}
% xurl：允许 URL 在任意字符处断行（中文书里长英文链接溢出边距的根治方案）
\usepackage{xurl}
% scheme=plain 下 ctexbook 不本地化目录名，显式改中文
\renewcommand{\contentsname}{目录}

% =========================================================
% 章扉页样式（重定义 \chapter）
% =========================================================
% chapter 显示为「第 N 章　标题」（同一行，章号朱砂色 + 标题黑体大字）
\ctexset{
  chapter = {
    beforeskip = 0pt,
    afterskip  = 40pt,
    pagestyle  = empty,
    lofskip    = 0pt,
    lotskip    = 0pt,
    name       = {第,章},
    number     = \arabic{chapter},
    nameformat = {\color{zhusha}\bfseries},
    titleformat= {\Huge\heiti},
    format     = {\centering},
    aftername  = \hspace{1em},
    aftertitle = {\par\vspace{18pt}\noindent\textcolor{zhusha}{\rule{\textwidth}{0.6pt}}\par},
  },
  section = {
    format     = {\large\bfseries\heiti\raggedright\color{mohei}},
    aftername  = \hspace{0.5em},
    beforeskip = 18pt,
    afterskip  = 8pt,
  },
  subsection = {
    format     = {\normalsize\bfseries\heiti\raggedright\color{mohei}},
    aftername  = \hspace{0.5em},
    beforeskip = 12pt,
    afterskip  = 6pt,
  },
}

% 目录：只显示到 section（N.M），不显示 subsection（N.M.K），避免目录过密
\setcounter{tocdepth}{1}

% 让目录里的章节编号（如 3.1、10.4.1）盒子宽一点，避免编号和标题分行
\makeatletter
% numberline 在编号后留一点空格，让「第N章　标题」编号和标题不挤在一起
\renewcommand{\numberline}[1]{%
  \hb@xt@\@tempdima{#1\hfil}%
  \hspace{0.5em}%
}
% 调整 section 在目录里的整体缩进和编号盒子宽度
\renewcommand*\l@section{\@dottedtocline{1}{1.5em}{2.8em}}
% chapter 在目录里给「第1章」编号盒子更宽（默认不够，会换行）
\renewcommand*\l@chapter{\@dottedtocline{0}{0em}{4.5em}}
\makeatother

% 黑体：ctex 已内建 \heiti，绑定到 sans CJK 族（上面 \setCJKsansfont）
% 无需重复定义

% =========================================================
% 引文框环境（epigraph）：朱砂左竖线 + 楷体小字 + 浅米灰底
% 用 tcolorbox（比 mdframed 更稳定）
% =========================================================
\newtcolorbox{epigraph}{
  enhanced,
  colback=qianhui,
  colframe=zhusha,
  boxrule=0pt,
  leftrule=2.5pt,
  arc=0pt,
  outer arc=0pt,
  left=14pt, right=14pt, top=10pt, bottom=10pt,
  before skip=14pt, after skip=14pt,
  fontupper={\kai\small\color{mohei}},
}
% =========================================================
% 页眉页脚
% =========================================================
\pagestyle{fancy}
\fancyhf{}
\renewcommand{\headrulewidth}{0.4pt}
\renewcommand{\footrulewidth}{0pt}
\fancyhead[LE]{\small\color{shenhui}{{BOOK_TITLE}}}
\fancyhead[RO]{\small\color{shenhui}\leftmark}
\fancyfoot[C]{\small\color{shenhui}\thepage}

% 章扉页用 empty
\fancypagestyle{plain}{\fancyhf{}\fancyfoot[C]{\small\color{shenhui}\thepage}\renewcommand{\headrulewidth}{0pt}}
